% 本文件由 LaTeX 排版 Agent 根据有效 translation.json 自动生成。
% author=Dionysius of Rome; work_id=0713; title=驳斥撒伯流派

\chapter{驳斥撒伯流派}

% structure: p0001 source=h1 target=omitted reason=page-title-already-rendered-by-parent

1. 如今真正应该与那些人争辩，他们通过分割和撕裂天主教会最庄严的宣告——独一统治——几乎如同分为三个权能、三个不同的实体（\OriginalTerm{实体}{hypostases}）和三个神祇，从而摧毁了它。因为我听说，你们中间有些宣讲和教导天主圣言的人，正是这种意见的老师；他们与撒伯流的意见几乎可以说是截然相反。因为撒伯流亵渎地说，圣子就是圣父，反之亦然；而这些人则是在某种意义上宣告三个神，因为他们将神圣的合一分割为三个彼此完全分离的不同实体。因为神圣的圣言必须与万有的天主结合，圣神也必须居住并停留在天主内；因此，神圣的天主圣三应当被归结并聚集到一内，如同到一个首脑内——即万有的全能天主。至于愚昧的马尔西翁的教导，他将独一统治割裂并分割为三个元素，这无疑是出自魔鬼，而不是基督的真正门徒或那些接受救主教导之人的。因为这些人确实正确知道，神圣圣经中宣告了天主圣三，但关于有三个神的教导，在旧约和新约中都没有教导。\par

2. 但那些认为圣子是受造之物，并断定主如同那些真正受造之物一样被造成的人，也并非少受责备；而神圣的宣言却见证祂是生成的，是适宜和恰当的，而非被创造或被造成的。因此，说主在任何方面是被手造的，这不是轻微的，而是极大的不敬。因为如果圣子是被造的，就有一段时间祂不存在；但祂永远存在，因为正如祂自己所宣告的，祂无疑是在圣父内。如果基督是圣言、智慧和德能——因为神圣的著作告诉我们基督就是这些，正如你们自己所知的——那么这些无疑是天主的德能。因此，如果圣子是被造的，就有一段时间这些德能不存在；这样，就有一段时间天主没有这些，这是完全荒谬的。但我为什么要向你们更详尽地谈论这些事呢？因为你们是充满圣神的人，尤其明白认为圣子是被造的论点会导致何等荒谬的后果。在我看来，这种观点的倡导者对这些事很少留意，因此完全偏离了正轨，以一种与神圣和先知圣经的要求相悖的方式来解释经文。\Quote{主创造了我在祂道路的起始。}因为，如你们所知，\Quote{创造}一词有不止一种含义；在这里，\Quote{创造}等同于\Quote{派定}在祂自己所造的工作之上——我是指由圣子自己所造的工作。但这个\Quote{创造}不应理解为与\Quote{制造}相同。因为制造与创造是不同的。\Quote{难道祂自己不是你的父亲，拥有你并创造了你吗？}梅瑟在申命记的伟大歌中如此说。因此，任何人都可以合理地驳斥这些人。哦，鲁莽和急躁的人！那么，\Quote{一切受造物的首生者}是受造之物吗？——\Quote{在晨星之前从胎中生出的那一位？}——祂以智慧的位格说：\Quote{在一切山丘之前，祂生了我。}（\BibleRef{箴 8:25}）最后，任何人可以在神圣话语的许多地方读到，圣子被称为是生成的，但从未被称为是被造的。从这些考虑，那些敢于说祂神圣而不可言喻的生成是被造的人，显然被证明在关于主的生成上思考了虚假的事。\par

3. 因此，那令人赞叹的神圣合一既不应被分割为三个神祇，也不应因将\Quote{受造}之名加于主而贬低主的尊贵和卓越的伟大，而应相信全能的天主圣父、祂的圣子耶稣基督，以及圣神。此外，圣言与万有的天主结合，因为祂说：\Quote{我与父原是一体。}（\BibleRef{若 10:30}）和\Quote{我在父内，父在我内。}（\BibleRef{若 14:10}）这样，无疑将保持神圣天主圣三的完整教义和独一统治的神圣宣告。\par

\SourceRecord{Translated by Robert Ernest Wallis.；Ante-Nicene Fathers；Vol. 7.；Edited by Alexander Roberts, James Donaldson, and A. Cleveland Coxe.；Buffalo, NY: Christian Literature Publishing Co.,；1886.；<http://www.newadvent.org/fathers/0713.htm>.}
